\section{REINFORCE}
    \subsection{Main Idea}
        \textbf{REINFORCE} is the most basic form of policy gradient method.
        It's the vanilla policy gradient method.
    \subsection{Objective Function}
        REINFORCE uses following objective function;
        \begin{equation}
            J(\pi_\param)=\E_{\tau \sim \pi_\param}[G_0] = \E_{\tau \sim \pi_\param}\left[ \sum_{t=0}^{T-1}\gamma^tr_{t+1} \right].
        \end{equation}
        We use some tricks to compute the gradient of $J(\pi_\param)$
        % TODO: add reference to trick here.
        Then we get
        \begin{equation}
            \nabla_\param J(\pi_\param) = \E_{\tau \sim \pi_\param}\left[ \sum_{t=0}^{T-1}G_t\nabla_\param\log{\pi_\param(a_t|s_t)} \right].
        \end{equation}
        Now we can use this to update the parameter $\param$. Since we cannot compute the exact expectation, we use Monte Carlo sampling
        to approximate $\nabla_\param J(\pi_\param)$ as following:
        \begin{equation}
            \nabla_\param J(\pi_\param) \approx \sum_{t=0}^{T-1}G_t\nabla_\param\log{\pi_\param(a_t|s_t)}
        \end{equation}
    \subsection{Pseudocode}
        \Algorithm{reinforce}
        {REINFORCE}
        {
            \inputitem{$\param$}{initial parameter} \\
            \inputitem{$\gamma$}{discount rate} \\
            \inputitem{$\alpha$}{learning rate}
        }
        {
            \inputitem{$\pi_\param$}{learned policy}
        }
        {
            initialize policy $\pi_\param$\;
            \Repeat{convergence}
            {
                generate episode $\{s_0,a_0,r_1,\dots,s_{T-1},a_{T-1},r_{T}\}$ following $\pi_\param$\;
                \For{$t=0,1,\dots,T-1$}
                {
                    \tcp{compute return-to-go}
                    $G_t=\sum_{k=t+1}^{T} \gamma^{k-t-1}r_k$\;
                    \tcp{policy update}
                    $\param \leftarrow \param + \alpha\nabla_\param\log{\pi_\param(a_t|s_t)}G_t$\;
                }
            }
            \Return{$\pi_\param$}
        }
    \subsection{Implementation Notes}
        \begin{itemize}
            \item When computing return-to-go, we can iterate backward from the terminal time to start, since $G_t \leftarrow \gamma G_{t+1} + r_{t+1}$ where $t=0, 1, 2,\dots,T-1$.
        \end{itemize}
        % TODO: add Implementation details
    \subsection{Pros and Cons}
        Pros:
        \begin{itemize}
            \item No need for model;
            \item Can cover both discrete and continuous actions;
        \end{itemize}
        Cons:
        \begin{itemize}
            \item Large variance;
            \item May converge to suboptimal policy;
        \end{itemize}